%==============================================================================
%  Chapter 5 — Discussion
%
%  Each section ties one empirical finding cluster (typically one of the
%  six confirmed Pass 2 patterns) to one theoretical construct, engages
%  with the most credible rival reading, and lands the interpretive move
%  for that cluster. Patterns are picked up in the order recommended by
%  synthesis_v2.md §6.2, with cross-pattern integrations as flagged in
%  §6.3. A new section (sec:disc-equity) is added for Pattern C
%  interpretation, which the original skeleton omitted.
%==============================================================================
\chapter{Discussion}
\label{ch:discussion}

\section*{Chapter overview}

The chapter interprets the findings of \Cref{ch:results} through the
constructs introduced in \Cref{ch:background}. Each section ties one
finding cluster to one theoretical construct, develops the
interpretation, and engages with the most credible rival reading. The
chapter then sets out two cross-cutting rival explanations of the
corpus as a whole, articulates four theoretical contributions, draws
practical implications for three audiences with a stake in the
instruments, and closes with the scope conditions under which the
findings travel.

%==============================================================================
\section{Logic succession as institutional design}
\label{sec:disc-succession}

The Science~$\to$~Market~$\to$~State succession across the innovation
pipeline (\cref{sec:res-profile}) and its intra-chapter reproduction
within every Challenge chapter (Pattern~D, \cref{sec:res-pattern-d})
together support an architectural reading of the Work Programme. The
instrument is logic-zoned by pipeline position, not by what applicants
do. Each instrument enacts a different configuration of
\textcite{thornton2012institutional}'s four institutional logics, and
the configuration is fixed by where in the pipeline the instrument
sits before any applicant arrives.

The macro-level succession is easy to mistake for an obvious
consequence of TRL placement: low-TRL science instruments foreground
science, high-TRL commercialisation instruments foreground markets.
Pattern~D weakens that reading. The Challenge chapters reproduce the
same four-step succession (state scope, science objectives,
professional governance, state impacts) across energy materials,
biotechnology, and artificial intelligence research. Substantively
different technology domains receive the same logic order. The
repeated configuration is therefore a feature of the chapter template
rather than a derivation from the science. Thornton's
interinstitutional system is not configured locally by each Challenge;
it is templated centrally and applied.

This finding extends the institutional-logics literature in two
directions. First, it locates institutional reproduction at the
chapter-template level: a meso scale between the field-level coercive,
mimetic, and normative mechanisms of \textcite{dimaggio1983iron} and
the micro-level coding of individual passages. The reproduction does
not require any individual sub-section to be coded twice or to
acknowledge the others; the template carries the order. Second, the
contrast between universal templating (Pattern~C, gender mainstreaming)
and adaptive entry (sustainability's four mechanisms,
\cref{sec:res-sustainability}) shows that the architecture distributes
different kinds of state priority through different channels.
Cross-cutting normative commitments travel as universal inserts;
substantive policy content travels through subject-matter-appropriate
channels. Both routes are state-organised.

The rival reading is that pipeline position is exogenous and the logic
order follows from it mechanically. The pipeline reading is not wrong;
it is incomplete. It cannot account for the Pattern~A embeddings that
put state-defined Challenge relevance inside the Excellence criterion
([082]) or for the market-readiness requirements at the science-phase
Specific Objectives ([102], [109]). The architecture is not only
ordered by pipeline position; it is also actively configured to
cross-embed logics where the pipeline would predict separation. The
cross-embeddings are where the rest of the chapter focuses.

%==============================================================================
\section{Tension-dissolution as legitimacy work}
\label{sec:disc-tension}

Pattern~A is best read through Suchman's typology of legitimacy as a
cognitive-legitimacy strategy. Presenting two logics as naturally
complementary does institutional work that an argument for their
compatibility would have to perform discursively. Where an argument
invites disagreement, complementarity by syntactic juxtaposition
foregoes the disagreement entirely. The instrument is presented as if
its internal contradictions did not need to be named.

The cleanest case is [082], where state-defined Challenge relevance is
embedded as the primary sub-criterion of scientific excellence. The
text does not argue that scientific excellence and state strategic
relevance are compatible; it lists the two as parallel quality
dimensions of the same applicant proposal. An applicant who reads the
criterion as a science criterion is being told to think of state
strategic priorities as quality of science. The configuration is not
deceptive (the wording is open) but it does sidestep the debate the
configuration could provoke.

The strategy is recognisable from the hybridisation literature on
institutional contradictions. \textcite{friedland1991bringing} treat
contradictions between institutional orders as the engine of
institutional change: actors expose contradictions in order to
mobilise alternative logics. Pattern~A is the inverse move. The Work
Programme assembles configurations that the literature would treat as
contradictory and presents them as continuous. The result is a text in
which structural conflict is not absent but rhetorically suppressed;
the readthrough surfaces the conflicts the text passes over.

Pattern~A is the most consequential cognitive-legitimacy strategy in
the corpus because it operates pre-discursively.
\textcite{scott2014institutions}'s cognitive pillar's diagnostic move
is to ask which taken-for-grantedness an institution recruits.
Pattern~A recruits the taken-for-grantedness of bureaucratic prose
itself: lists, parallel sub-criteria, and conjunction. The
\enquote{and} that joins science to market is doing more work than its
grammatical role would suggest.

%==============================================================================
\section{Coercive governance under an autonomy register}
\label{sec:disc-coercive}

Pattern~B is best read as coercive isomorphism through grant-agreement
mechanics, with the implication that
\textcite{dimaggio1983iron}'s coercive mechanism extends in the EIC
case beyond legislative source to contractual discretionary authority.
The EIC formally protects researcher autonomy at [021] while retaining
hard override authority over capital continuity at [024] and proactive
roadmap-setting authority at [026]. The gap between stated autonomy
and operative control is the substantive observation; the
institutional reading is that the gap is held open by the legal form
of the grant agreement.

The autonomy promise is regulatively underwritten only at its outer
edge. Applicants choose their research direction within the scope of
their proposal; the proposal is judged against a Challenge whose
relevance the state defines; the project is overseen by a Programme
Manager who can suspend or terminate the grant where the project no
longer addresses that Challenge. The autonomy is real at the level of
methodological choice. It is sharply bounded at the level of
direction. The Programme Manager is a hired technocrat, neither a
peer scientist nor an elected official, and the directive authority
runs through the legally neutral language of grant administration.

The most credible rival reading is that Pattern~B is
\texttt{LOGIC\_STATE} operating through professional instruments
rather than a distinct administrative-managerial logic, and that the
\texttt{ISO\_COERCIVE} extension introduced by coding rule~3 risks
collapsing useful analytical distinctions. The rival is right that no
new logic needs to be posited; this thesis does not propose one. The
rival is wrong about the coercive extension.
\textcite{scott2014institutions}'s regulative pillar treats the grant
agreement as a legal instrument just as it treats a regulation, and
the suspension and termination clauses operate through that
instrument. Coercive isomorphism in the EIC therefore does not require
a legislative source on each occasion: the legislative source
authorises the contractual form, and the contractual form does the
proximal coercive work.

The institutional design feature this exposes is that the EIC mediates
between political principal and scientific applicant through a hired
official whose discretion is contractually authorised, not
professionally constrained. Whether the mediation should be read as a
deliberate design choice (state oversight wrapped in administrative
discretion) or as a constitutive contradiction (the autonomy promise
and the termination clause cannot both fully hold) depends on how one
reads the silence at [061], where the dual-responsiveness mandate is
asserted without specifying how competing directions are reconciled.
The first reading is more generous to the instrument; the second is
more honest about its structural tension.

%==============================================================================
\section{TRL as cross-pillar legitimation}
\label{sec:disc-trl}

Pattern~F is the analytically densest single site in the corpus
because the Technology Readiness Level framework carries three
distinct legitimation axes simultaneously. As \texttt{REG\_TRL} the
framework is regulative: it sets eligibility thresholds, draws scope
boundaries, and routes budget along the pipeline. As
\texttt{COG\_LINEAR} it is cognitive: a developmental ladder so
taken-for-granted in EU innovation policy that its use in the Work
Programme reads as scientific measurement rather than as policy
choice. The two pillars co-operate within the same passages; the
cognitive surface provides scientific-register cover for the
regulative gate beneath.

The third axis, identified at the Pass~2 mimetic finding, is the
framework's US-defence provenance. TRL originates in NASA technology
management and was adopted across US-defence procurement before its
uptake in EU research policy. The provenance carries mimetic
legitimacy in \textcite{dimaggio1983iron}'s sense: the EU borrows the
framework partly because it has been seen to work in a reference
field. The borrowing is not exposed by the Work Programme; the TRL
ladder is presented as a neutral developmental schema rather than as
an inherited artefact.

Read jointly, the three axes do a specific kind of work. The political
question of where in the innovation pipeline public money should flow
is reframed as a technical question of developmental readiness.
\textcite{scott2014institutions}'s cognitive pillar supplies the
scientific register; the regulative pillar performs the gatekeeping;
mimetic isomorphism imports the framework from a field in which it
carries authority. The result is an artefact whose political
consequences (which research is fundable, which is too early, which is
too late) are routinely presented as technical classifications.

The cross-link to Pattern~A is that both are cognitive-legitimacy
strategies, but they operate at different scales. Pattern~A is a
sentence-level move that joins logics without argument; Pattern~F is
an artefact-level move that does the same work through a designed
schema. The implication for Suchman's framework is that cognitive
legitimacy in the EIC operates not only through discourse but through
designed artefacts. The TRL ladder and (as the next section but one
develops) the Sovereignty Seal are objects, not rhetorical features,
and their design carries the legitimation strategy.

%==============================================================================
\section{The state social-equity template as developmental-state
insertion}
\label{sec:disc-equity}

Pattern~C is best read as a second developmental-state dimension of
the EIC, parallel to but distinct from the catalytic-investment
dimension developed in \cref{sec:disc-step}. The pattern shows the
state inserting its own social-equity priorities (gender balance,
female-researcher empowerment, underrepresented-group inclusion) into
an innovation funding instrument and ranking those priorities above
the instrument's own market and geographic priorities in the
tiebreaker architecture at [059]. The interpretive weight is on that
ranking.

An innovation instrument whose stated success criteria are
technological excellence, market creation, and strategic-technology
development might be expected to break ties on dimensions internal to
those criteria: SME participation, geographic cohesion across Member
States, contribution to a critical technology area. The tiebreaker
does the opposite. It ranks gender balance first, female-researcher
participation second, and so on, with SME and geographic
considerations below. The instrument is therefore being used to carry
state-level political commitments whose direct relationship to the
instrument's stated success criteria is not asserted.

Read through \textcite{mazzucato2018mission}'s mission-oriented
investment frame, this is a recognisable developmental-state move: the
state pursues collective political interests through the investment
instrument, not despite it. The catalytic-investment mechanism
explored at Pattern~E is the financial dimension of the same approach;
Pattern~C is its social-policy dimension. Both operate through the
same instrument; both are legitimised by different vocabularies
(catalytic-investment language for Pattern~E, gender-mainstreaming
language for Pattern~C); and both put state priorities at the centre
of an instrument whose surface presentation is technological.

The standardised wording of the template (near-identical at [047] and
[063]) suggests that the insertion happens at the legislative layer
rather than at the EIC's own programmatic layer. The Horizon Europe
Regulation imposes gender-mainstreaming obligations on the whole
programme, and the Pattern~C passages reproduce those obligations as
compliance text. Two consequences follow. First, Pattern~C is coercive
isomorphism in the classical \textcite{dimaggio1983iron} sense: a
legislative mandate reproduced across instruments by the agencies
subject to it. Second, the tiebreaker ranking is the EIC's own move:
the legislative mandate requires gender mainstreaming, but the
ordering of priorities in the tiebreaker is a programmatic choice the
EIC made on top of the mandate. The mandate is coercive; the ranking
is the EIC's reading of how strongly to enforce it.

%==============================================================================
\section{STEP as a developmental-state instrument}
\label{sec:disc-step}

STEP is best read not as the EU adopting market-VC logic but as a
developmental-state instrument deploying market mechanics. The three
sub-sections develop the reading across the circular market validation
finding (Pattern~E), the Sovereignty Seal, and the explicit mimetic
benchmarking against DARPA, Temasek, and the Chinese Government
Guidance Funds. Taken together they support treating STEP as the EU's
clearest instance of mission-oriented state investment to date.

\subsection{Circular market validation as deliberate design}
\label{sec:disc-step-circular}

Pattern~E (\cref{sec:res-pattern-e}) is best read as a deliberate
developmental-state design rather than as a contradiction. The
simultaneous requirements of demonstrated market interest ([140],
[141]) and demonstrated market failure ([176]) are only satisfiable
through the EIC's prospective 20--33\% anchor participation ([145]).
Read as a contradiction, the configuration shows the state claiming to
test market readiness while eliminating the independent market
assessment it claims to test for. Read as a design, the configuration
shows the state constructing the market signal it then uses as
legitimation for its own intervention.

The second reading is more consistent with what STEP actually does.
The instrument is not an independent market evaluator; it is a
catalytic investor. \textcite{mazzucato2018mission}'s patient-capital
account treats the state as an anchor investor in technology
development that the market alone will not fund at the relevant scale
or risk profile. Chinese Government Guidance Funds are structured
around precisely the configuration Pattern~E describes: state-anchor
participation that crowds in private co-investment, with the
co-investment serving as both validation and risk-sharing mechanism.
The two are not independent. The state co-investment creates the
private commitment which the state then takes as evidence of demand.

Read in that lineage, the circularity is not a flaw of the instrument
but a feature of the class of instruments to which STEP belongs. What
is unusual in the EU case is the surface presentation. STEP narrates
itself in the vocabulary of market gap, deal flow, and investor
co-investment (the vocabulary of liberal-market venture finance) while
operating with a state-anchored structural design. The vocabulary is
consistent with \textcite{ferrary2009vc}'s description of VC
ecosystems; the structure is consistent with the Chinese GGF model.
The interpretive consequence is that STEP is best classified by its
structure rather than its vocabulary. The vocabulary is the register;
the structure is the logic.

\subsection{The Sovereignty Seal as field-bridging legitimacy}
\label{sec:disc-step-seal}

The Sovereignty Seal is best read as a designed legitimacy artefact in
Suchman's pragmatic-plus-moral sense, with the additional property
that it is engineered to travel across institutional fields. As a
pragmatic-legitimacy device, the Seal gives access to resources:
InvestEU Venture Debt, Business Acceleration Services, and a signal
usable in private fundraising. As a moral-legitimacy device, it
carries the state's certification of strategic alignment with EU
interests. The Suchman reading is straightforward at this level.

The analytically interesting property is the decoupling at [166]. The
Seal can be awarded on a NO~GO investment decision where merit
criteria are met but budget is unavailable. Stripped of the capital
commitment, the Seal is no longer a side-effect of an investment
decision; it is an instrument in its own right. The state certifies
strategic value independently of whether it commits capital, and the
certification is designed to be portable. It is a credential issued in
one institutional field (the EIC's award procedure) and intended to be
honoured in others (InvestEU, private investment, partnership
negotiation).

This is an institutional innovation worth naming. The literature on
state certification typically treats certification as a side-effect of
regulation or of capital provision. The Seal is neither. It is a
field-bridging legitimacy artefact: a designed object whose function
is to carry state-issued legitimacy from the field of award to fields
of subsequent use, independently of the capital decision that produced
it. The EU has not previously deployed a portable state certification
of this kind in innovation finance, and the design is worth tracking
across future instruments.

\subsection{Mimetic isomorphism against DARPA, Temasek, and Chinese
GGFs}
\label{sec:disc-step-mimesis}

The explicit benchmarking of STEP against DARPA, Temasek, and the
Chinese Government Guidance Funds is mimetic isomorphism at the
instrument-design level. \textcite{dimaggio1983iron}'s account treats
mimetic isomorphism as borrowing from organisations that appear
successful in conditions of uncertainty; STEP cites the three
referents directly and adopts elements of each. What STEP reproduces
is the equity and catalytic-investment mechanics, the state-anchor
position, and the orientation toward strategic technologies. What it
does not reproduce is the operational autonomy of the cited agencies
from their political principals. DARPA in particular is characterised
by an unusual degree of programme-officer autonomy from Congressional
priorities; STEP retains Commission Award Decisions, quarterly
evaluation batches, and the standard EU procurement architecture.

The hybrid that results sits closer to the Chinese GGF model than to
DARPA. The Chinese GGF structure is built around state-anchor
participation with private co-investment, operating through
state-controlled administrative architecture and oriented toward
state-defined strategic priorities. That description fits STEP more
cleanly than any DARPA-like reading does. The DARPA citation is
nonetheless analytically meaningful: it carries cognitive legitimacy
in EU policy debates that the GGF reference cannot. The DARPA name
does legitimation work; the GGF structure does the design work.

The cross-link to Pattern~F (\cref{sec:disc-trl}) is that mimetic
borrowing is not confined to STEP. The TRL ladder gates Pathfinder and
carries the same US-defence provenance. The EIC's institutional logic
is therefore partly constituted by mimetic borrowing from US
innovation finance and US-defence procurement, operating at two
different points in the pipeline. The borrowing is selective in both
cases. The reference fields supply legitimacy and design elements; the
EU procurement architecture absorbs them.

%==============================================================================
\section{Rival explanations}
\label{sec:disc-rival}

Two cross-cutting rival readings of the corpus deserve direct
engagement.

\paragraph{The realpolitik strategic-autonomy reading.}
The first reading is that the underlying logic of the Work Programme
is geopolitics rather than institutional theory. STEP exists because
the EU is in strategic competition with the US and China on critical
technologies; the equity and catalytic-investment mechanics are
selected because the strategic competition requires them; the
Sovereignty Seal is a sovereignty instrument. On this reading, the
institutional-theoretic apparatus is descriptively accurate but
explanatorily redundant.

The reading captures something real. The strategic-autonomy frame is
explicit in the corpus, particularly at the Strategic Goals level and
in the STEP introduction. The institutional reading does not displace
it; it sharpens it. The strategic-autonomy frame explains why the EU
has built a developmental-state-style instrument; it does not explain
why the instrument carries the particular features the corpus records
(the Pattern~A complementarity, the Pattern~B authority structure, the
Pattern~C tiebreaker, the Pattern~F artefact triple-load). Those
features are not derivable from the strategic-autonomy premise. They
are institutional choices made under that premise, and the
institutional reading is what gives the choices their texture.

\paragraph{The path-dependence / budget-cycle reading.}
The second reading is that STEP is the predictable next move in an EU
funding cycle. The EU funds programmes in seven-year frameworks,
mid-term reviews routinely produce new instruments, and STEP is the
mid-term-review outcome for the current cycle. On this reading the
corpus reflects what the budget cycle requires rather than any deeper
institutional construction.

Again the reading captures something real and is not displaced by the
institutional reading. The institutional reading explains the
configuration of STEP, not the timing of its appearance. The
configuration --- equity-only with state anchor, Sovereignty Seal as
portable artefact, DARPA and GGF mimetic citation --- is not forced by
the budget cycle. Other configurations were available, and the EU
chose this one. The institutional reading is what makes the choice
legible.

\paragraph{Net.}
The institutional reading does not compete with strategic-autonomy or
path-dependence accounts; it operates at a different level. The two
rivals explain why an instrument of this class exists at this moment.
The institutional reading explains why the instrument has the texture
it does. The thesis claims explanatory traction at the second of those
levels, not the first.

%==============================================================================
\section{Theoretical contributions}
\label{sec:disc-contrib}

The thesis offers four contributions to the institutional-theoretic
literature on supranational public-investment instruments.

\begin{enumerate}[leftmargin=*]
  \item \textbf{Extension of the institutional-logics framework to
    supranational public-investment instruments.} The
    \emph{logic-as-register} distinction (rule~7 in the Pass~2
    codebook) treats market vocabulary as analytically distinct from
    market organising logic. The distinction is generally useful for
    instruments in which the state borrows from market discourse
    without transferring authority to market actors. The EIC
    instruments are one such case; the distinction generalises beyond.
  \item \textbf{\emph{Tension-dissolution} as a distinct legitimacy
    strategy.} Pattern~A names the textual move that presents
    conflicting logics as naturally complementary, analytically
    distinct from surface-visible logic conflict (the
    \texttt{LOGIC\_TENSION} code) and from
    \textcite{friedland1991bringing}'s mobilisation of contradiction.
    The strategy operates pre-discursively through bureaucratic prose
    conventions, and it merits naming inside Suchman's
    cognitive-legitimacy typology.
  \item \textbf{The Sovereignty Seal as a portable field-bridging
    legitimacy artefact.} The Seal is a designed credential whose
    legitimacy is engineered to travel across institutional fields,
    decoupled from the capital decision that produced it. The artefact
    extends the literature on state certification beyond
    side-effect-of-regulation accounts, and is worth tracking across
    subsequent EU innovation instruments.
  \item \textbf{\emph{Circular market validation} as a structural
    property of developmental-state investment design.} Pattern~E is
    not a flaw of STEP; it is a recognisable feature of the class of
    instruments to which STEP belongs, observable in Chinese
    Government Guidance Funds and consistent with
    \textcite{mazzucato2018mission}'s patient-capital model. The label
    \emph{circular market validation} names the configuration of
    simultaneous market-interest and market-failure requirements
    satisfiable only through state-anchor participation.
\end{enumerate}

A fifth observation, more modest in scope, is that Pattern~F's
triple-load operates as a cross-pillar legitimation mechanism within a
single artefact. \textcite{scott2014institutions}'s three pillars are
typically treated as analytically separable mechanisms; TRL shows that
they can co-operate within a single designed object. The finding
refines rather than contradicts Scott's frame.

%==============================================================================
\section{Practical and policy implications}
\label{sec:disc-practice}

\subsection{For EIC programme design}
\label{sec:disc-prog-eic}

The actionable implications for the EIC are about what the instrument
should make explicit rather than about what it should change. The
developmental-state character of STEP is best stated plainly rather
than narrated through market vocabulary; the Sovereignty Seal would be
more useful to applicants and to downstream institutions if its
decoupling from the investment decision were stated explicitly in the
call text rather than buried at [166]; the Programme Manager's
directive authority is more defensible if it is presented as a design
feature rather than as a qualification to the autonomy promise at
[021]. The trade-offs are real. The market vocabulary protects STEP's
standing in EU policy debates that are hostile to industrial policy;
the autonomy language at [021] reflects a genuine commitment that the
termination clauses qualify but do not contradict. The recommendation
is for explicitness about the trade-offs, not for their resolution.

\subsection{For applicants and intermediaries}
\label{sec:disc-prog-app}

Applicants and the consultants and accelerators who advise them should
read each instrument's evaluation in terms of the logic that actually
organises the criterion, not the vocabulary in which it is expressed.
The Pathfinder Excellence criterion is a science criterion with
state-strategic-relevance embedded as its primary sub-criterion
([082]); proposals that ignore the embedding are at a disadvantage
that the criterion's name does not announce. The STEP market-interest
criterion is satisfied in practice through co-investment around the
EIC anchor; proposals that present independent market interest without
coordination to the EIC's prospective anchor position are working
against the architecture. The Sovereignty Seal is a credential worth
pursuing for its field-bridging properties, independently of whether
the underlying investment decision is favourable.

\subsection{For EU policy makers}
\label{sec:disc-prog-eu}

The findings have two implications for the design of mission-driven
funding instruments in FP10 and successor frameworks. First, the
developmental-state inflection of STEP is best treated as a template
rather than a one-off. The structural features that produce the
inflection --- state-anchor participation, portable certification, and
mimetic borrowing from reference fields --- are reusable across
mission-oriented instruments and are likely to recur. Second, the
tiebreaker architecture at [059] is the simplest available point at
which state-level political priorities are weighted against
instrument-internal performance criteria. Adjustments there would be
a more transparent way to recalibrate priorities than rewording the
substantive criteria. The recommendation is to treat the tiebreaker as
a deliberate design surface rather than as administrative
housekeeping.

%==============================================================================
\section{Boundary conditions}
\label{sec:disc-boundary}

The findings travel to other EU instruments with comparable
architectures (the Innovation Fund, and an EU Sovereignty Fund if
established) and to mission-oriented public-investment instruments in
coordinated-market economies more generally. The Pattern~F triple-load
reading travels wherever TRL is institutionally embedded in a
regulative gate function. The Pattern~E reading travels wherever a
state-anchor investment vehicle requires demonstrated market interest
at scales the anchor itself reaches. The Pattern~C reading travels
wherever a legislatively mandated social-policy template is reproduced
across instruments by the agencies subject to it.

The findings travel less well in two directions. First,
mission-oriented instruments in liberal-market national systems may
behave differently. The Pattern~A complementarity strategy presumes a
policy environment in which industrial-policy vocabulary is contested;
in environments where industrial policy is openly named, the
rhetorical work that Pattern~A performs is unnecessary, and the
pattern itself may not appear. Second, the readings are anchored in
the Pathfinder and STEP Scale Up sections of the EIC Work Programme
2026. They are not claims about the EIC Accelerator, which uses a
different evaluation architecture, nor about EIC grant instruments
outside the Challenge structure. The scope is the corpus defined in
\Cref{ch:methodology}, and claims should be read against that scope.
